% Modernised reading — fonds Alexandre Grothendieck, shelfmark 141 (pages 1-15).
%
% This is an INTERPRETATION, not a transcription. It re-reads the pages in
% today's notation and vocabulary, states the hypotheses the manuscript leaves
% implicit, and drops the critical apparatus so that the mathematics reads
% straight through. Everything the brackets and underlines carried moves into
% a footnote. It is held to being mathematically correct as it stands; where
% the manuscript is loose or mistaken the true statement appears and the
% footnote says what the page has. In any dispute of substance the
% transcription governs: whoever cites Grothendieck cites batch-01.fr.tex.
%
% Scope: a single document for the whole folder. The shelfmark holds one
% batch (pages 1-15), read whole before any of this was written. Pages 2, 4
% and 6 are typed administrative notices, dated November 1980, with nothing
% in his hand.
%
% The folder has no single argument. It holds four runs on four kinds of
% paper — Kummer theory (1, 3, 5), classifying spaces (7-8), cartographic
% groups (9-10), pi-sets and bitorsors (11-15) — and the spine section says
% so rather than inventing a thread. Two connections between runs are this
% reading's and are flagged as such: the group pi_1(X, Gamma, a) of page 14
% is the combinatorial model of the pi_1(X, pi) of page 8; and the involution
% epsilon of pages 11 and 14 is the Z/2-extension of page 13, whose coherence
% condition is the F mu = mu F of page 15.
%
% Notation fixed for the whole folder. Groups act on the left unless said
% otherwise; Ens(pi) is the category of left pi-sets; a (pi', pi)-biset is a
% set with a left pi'-action and a commuting right pi-action, written
% {}_{pi'}X_{pi}; X x_pi E is the contracted product (quotient of X x E by
% (x lambda, e) ~ (x, lambda e)). Bit(pi', pi) is the groupoid of
% (pi', pi)-bitorsors, Isomext(pi', pi) = Isom(pi', pi)/Int(pi), Out(pi) =
% Isomext(pi, pi). H^1(X, pi) is the set of isomorphism classes of principal
% pi-coverings. mu_n, G_m, Pic as usual; the page's Rev(S, mu_n) is written
% H^1(S, mu_n).
%
% Substantive departures from the page, with pages.
%  - Page 1: « K*/K*^n = Z/nZ » and « every extension is cyclic » use that V
%    is henselian and that 1 + m and k* are n-divisible; supplied. The
%    residue-characteristic-p statement needs V henselian (the page only says
%    « pas nécessairement complet ») and is a statement about the TAME
%    quotient of pi_1(K); the full pi_1(K) has a pro-p wild inertia.
%  - Page 3: the map pi_1(S, sbar) -> lim mu_n(S) is written without saying
%    what defines it; it is read as the Kummer character of a unit a, and the
%    page's condition on n (read « premier à s ») as n invertible on S.
%  - Page 5: the first line omits the second H^0(S, G_m) and the n-th power
%    map; exactness of the Kummer sequence needs n invertible on S (étale
%    topology) or the fppf topology; a must be a unit.
%  - Page 7: the theorem needs X to range over CW complexes (or paracompact
%    spaces) and « contractile » is « faiblement contractile » in general.
%  - Page 8: « défini à isom. unique près » holds for the pair (B, class of E),
%    not for B alone, whose self-equivalences realise Out(pi).
%  - Page 9: the page never defines Gamma_{p,q} nor says which rotation p and
%    q bound; its two mentions of rho_s (rho_s^3 = 1 in Gamma_{infty,3},
%    rho_s^p alone) are compatible only with different conventions. Only
%    statements independent of that choice are made. « C_{p',q'} = C_{p,q}/R »
%    is restated as a full embedding with a left adjoint E |-> E/R.
%  - Page 9: the special value of the modular invariant is 1728 = 2^6.3^3; the
%    page has what reads 2^3.3^4.
%  - Page 11: 𝒞 must be cocomplete (read « avec lim-> »); the rule rho_1 |->
%    rho_infty rho_0 defines an automorphism only for the relation
%    rho_infty rho_1 rho_0 = 1 (equivalently rho_1 rho_0 rho_infty = 1), not
%    for rho_0 rho_1 rho_infty = 1.
%  - Page 13: the struck definition pi_I = L_I/(lambda_0...lambda_{n-1}) is
%    the ORIENTED group pi_I^+, which the page then writes; kept apart.
%  - Page 14: a first product law, on the left, does not compose; the law on
%    the right does and is the one given.
%  - Page 15: Bit(pi', pi) -> Isomext(pi', pi) is a bijection on isomorphism
%    classes; the groupoid also has automorphisms, the centre Z(pi), which the
%    page does not mention.
%
% Announced and never established in the folder: the equivalence 1) <=> 2) of
% Thm 1 (the card breaks off at « 2) E est contractile »); Thm 2 (stated, not
% proved); the equivalences Bit.invol(pi) = Invol(Ens(pi)) = Ext(Z/2, pi) of
% page 13 (stated; the reading sketches why they hold); the « carte trouée »
% of page 10, whose definition is left at its first line.
%
% The archivists' title « note Alexandre (bi tenseur…) » is not discussed by
% the pages; « bitorseur », the word of pages 12-15, is offered in the spine
% section as a likely reading of « bi tenseur », and only as that.
%
% Pass: Opus 5 (claude-opus-5), 2026-09-13 — first pass, unchecked against the pages by a human.
\documentclass[11pt,a4paper]{article}
% Le préambule commun à toutes les transcriptions du fonds Grothendieck.
%
% Il tient en une idée : **l'incertitude se compose, elle ne se résout pas.**
% Une transcription qui lisse un mot douteux a détruit la seule chose pour
% laquelle on la faisait — un lecteur doit pouvoir savoir, à chaque endroit,
% ce qui était sur le feuillet et ce qui a été ajouté. Les six macros qui
% suivent sont donc l'appareil critique, et non de la décoration.
%
% Les mêmes commandes sont reconnues par `scripts/render.mjs`, qui produit la
% vue de lecture du site. Une macro ajoutée ici doit l'être aussi là-bas,
% faute de quoi le rendu échouera bruyamment — ce qui est voulu : un rendu
% silencieusement faux est pire qu'un rendu absent.

\ProvidesPackage{grothendieck}[2026/08/08 Transcriptions du fonds Grothendieck]

\usepackage{amsmath,amssymb,amsthm}
\usepackage{mathrsfs}
\usepackage{stmaryrd}
% Les diagrammes commutatifs sont partout dans ce fonds : c'est la langue
% même de la géométrie algébrique de Grothendieck, et une transcription qui
% les rendrait en prose ne transcrirait rien.
\usepackage{tikz-cd}
% Français par défaut — c'est la langue des feuillets — mais l'anglais est
% chargé aussi : la traduction et le résumé sont des documents anglais, et la
% typographie française (espaces avant les deux-points) y serait une faute.
% Ces documents basculent par \selectlanguage{english} en tête de corps.
\usepackage[english,french]{babel}
\usepackage{fontspec}
\usepackage{geometry}
\geometry{a4paper,margin=25mm,marginparwidth=28mm}
\usepackage{marginnote}
\usepackage{xcolor}
% ulem, not soul. `soul`'s \st and \ul are built on a character-by-character
% scan that XeLaTeX's font machinery breaks: the document fails with an
% undefined control sequence rather than degrading. ulem does the same two jobs
% and is Unicode-engine safe. `normalem` stops it hijacking \emph, which the
% transcriptions use for Grothendieck's own emphasis.
\usepackage[normalem]{ulem}
\usepackage{hyperref}
\hypersetup{colorlinks=true,linkcolor=black,urlcolor=black,pdfborder={0 0 0}}

% --- Métadonnées du lot -----------------------------------------------------
% Reprises telles quelles par le script de rendu, qui les affiche en tête de
% la vue de lecture. Elles fixent ce que le fichier prétend couvrir : sans
% elles, une transcription flotte sans point d'attache dans un fonds de seize
% mille feuillets.

\newcommand{\folder}[1]{\gdef\grothFolder{#1}}
\newcommand{\batch}[1]{\gdef\grothBatch{#1}}
\newcommand{\leaves}[2]{\gdef\grothFirst{#1}\gdef\grothLast{#2}}
\newcommand{\foldertitle}[1]{\gdef\grothFoldertitle{#1}}
\gdef\grothFolder{?}\gdef\grothBatch{?}\gdef\grothFirst{?}\gdef\grothLast{?}\gdef\grothFoldertitle{}

% --- Le résumé --------------------------------------------------------------
%
% Il ouvre la lecture modernisée, sous le titre « Résumé », et il est composé
% en petit. Ce n'est pas un ornement : c'est le seul passage du document qui
% ne soit pas des mathématiques, et un lecteur doit pouvoir le reconnaître
% avant de l'avoir lu — puis le sauter s'il connaît déjà le sujet, ou s'y
% arrêter s'il ne le connaît pas. Le filet qui le suit marque la reprise du
% corps.

\newenvironment{resume}
  {\small\setlength{\parskip}{0.45em}}
  {\par\medskip\hrule\medskip}

% --- Mots-clés ---------------------------------------------------------------
%
% En anglais, en clôture du résumé de la lecture modernisée : le vocabulaire
% moderne sous lequel chercher ce que les feuillets construisent. Le manifeste
% du site (`npm run manifest`) les extrait pour étiqueter la cote entière —
% cette ligne est la source unique des tags, il n'y a pas de fichier à côté.
\newcommand{\keywords}[1]{\par\smallskip\noindent
  {\footnotesize\textsc{Keywords} --- \textit{#1}}\par}

% --- L'appareil critique ----------------------------------------------------

% Le numéro de feuillet, en marge. C'est l'ancre qui permet de régler un
% désaccord sur une formule en regardant *un* feuillet plutôt qu'une chemise
% de six cents. Le site s'en sert aussi pour tourner les pages du fac-similé
% quand on fait défiler la transcription.
\newcommand{\leaf}[1]{%
  \marginnote{\footnotesize\color{gray!70}\textsf{#1}}%
  \label{leaf:#1}\ignorespaces}

% Illisible : on ne devine pas. Un mot inventé qui se lit comme les autres est
% le pire résultat possible.
\newcommand{\ill}{\textcolor{red!60!black}{[\,\ldots\,]}}

% Lecture douteuse : le mot est donné, et signalé comme incertain.
\newcommand{\uncertain}[1]{\uline{#1}}

% Ajout de l'éditeur — une lettre manquante, un mot sous-entendu.
\newcommand{\add}[1]{\textcolor{blue!50!black}{[#1]}}

% Biffé par Grothendieck lui-même. Ce qu'il a rayé fait partie du document :
% une rature dit souvent où la pensée a changé de direction.
\newcommand{\struck}[1]{\sout{#1}}

% Note du transcripteur, sur l'état matériel du feuillet.
\newcommand{\note}[1]{\par\smallskip{\small\color{gray!60!black}\textit{[#1]}}\par\smallskip}

% Note marginale de Grothendieck, distincte du corps du texte.
\newcommand{\marginal}[1]{\par\smallskip\noindent\hspace{1em}%
  {\small\color{gray!60!black}#1}\par\smallskip}

% --- Titre ------------------------------------------------------------------

\AtBeginDocument{%
  \begin{center}
    {\large\bfseries Fonds Alexandre Grothendieck --- cote n\textsuperscript{o}~\grothFolder}\\[2pt]
    {\itshape\grothFoldertitle}\\[4pt]
    {\small feuillets \grothFirst--\grothLast\ (lot \grothBatch)}\\[2pt]
    {\footnotesize\color{gray!60!black}Transcription établie d'après le fac-similé numérisé
     par l'Université de Montpellier.\\
     Les lectures incertaines sont soulignées, les illisibles notées [\,\ldots\,],
     les ajouts entre crochets.}
  \end{center}
  \vspace{1em}}

\endinput


\folder{141}
\pages{1}{15}
\dating{[à partir de 1980]}
\watermark{Édition de démonstration\\interprétation personnelle de l'œuvre}
\foldertitle{Théorie de Kumer ? + note Alexandre (bi tenseur…) : notes manuscrites (s.d.) — lecture modernisée du dossier entier}

\begin{document}

\section*{Résumé}

\begin{resume}
Ces quinze pages tournent autour d'une seule question, prise sous quatre
angles : quand un objet se « déplie » en plusieurs copies de lui-même qu'un
groupe permute, qu'est-ce qui classe toutes les façons de le faire ?

L'image la plus simple est celle des racines. Au-dessus d'un nombre $a$, les
$n$ racines $n$-ièmes de $a$ forment un petit ensemble que le groupe des
racines de l'unité fait tourner sur lui-même ; faire varier $a$, c'est faire
varier ce « revêtement ». La \emph{théorie de Kummer}, à laquelle s'ouvrent les
pages 1 à 5, dit que sur certains corps toutes les extensions sont de cette
forme, et en tire une description complète de leur groupe de Galois ; puis
elle passe à un espace de base quelconque, où les revêtements de ce type sont
comptés par une suite exacte qui mêle les fonctions inversibles et les fibrés
en droites.

Les pages 7 et 8 font la même chose en topologie. Pour un groupe $\pi$, il
existe un espace $B$ — l'\emph{espace classifiant} — tel que se donner un
revêtement de groupe $\pi$ au-dessus d'un espace $X$ revienne exactement à se
donner une application de $X$ dans $B$, à déformation près : comme une carte
routière contient d'avance tous les trajets qu'on pourra y tracer. Ce qui
caractérise $B$ est d'une simplicité frappante : son grand revêtement doit être
contractile.

Les pages 9 et 10 sont des esquisses d'un tout autre ton, au feutre, pleines de
sphères, de polygones et de flèches. Ce sont les \emph{cartes} : des graphes
dessinés sur des surfaces, que Grothendieck décrit au même moment par un
groupe, le « groupe cartographique », agissant sur les drapeaux (sommet, arête,
face) de la carte. Une carte devient un ensemble muni d'une action de groupe, et
le groupe modulaire $SL(2, \mathbb{Z})$ apparaît au passage, comme le groupe
des triangulations.

Les pages 11 à 15, enfin, reprennent la question dans le langage le plus nu :
celui des ensembles sur lesquels un groupe agit. Quand deux groupes ont-ils
« les mêmes » ensembles à action ? Exactement quand il existe entre eux un
\emph{bitorseur} — un ensemble sur lequel l'un agit à gauche, l'autre à droite,
chacun librement et transitivement. Ce n'est pas tout à fait la même chose
qu'un isomorphisme : c'est un isomorphisme connu seulement à conjugaison près,
et cette nuance est le cœur des dernières pages, qui s'achèvent sur les
bitorseurs « involutifs », où revient, semble-t-il, la conjugaison complexe
agissant sur la sphère privée de quelques points.

Les feuillets datent vraisemblablement de 1980 ou peu après : les versos
administratifs sont de novembre 1980, quelques mois avant \emph{La Longue
Marche à travers la théorie de Galois} (1981), dont le dossier est voisin dans
l'inventaire, et quelques années avant l'\emph{Esquisse d'un programme} (1984),
où les cartes et le groupe cartographique sont exposés.

\keywords{Kummer theory, tame fundamental group, Kummer sequence, classifying space, Eilenberg–MacLane space, Borel construction, cartographic group, dessins d'enfants, extended modular group, free cocompletion, bitorsor, outer automorphism group, group extension, reflection group of an ideal polygon}
\end{resume}

\pagerange{1}{15}
\section*{Le fil du dossier, et les conventions}

Le dossier n'est pas un argument suivi. Il réunit quatre ensembles de
feuillets, qui se distinguent à l'œil — encre, papier, format — autant que par
le sujet ; aucun ne porte de titre de la main de Grothendieck.

\begin{itemize}
  \item \textbf{Pages 1, 3, 5 : théorie de Kummer.} Sur un anneau de
    valuation discrète, puis sur un schéma quelconque ; stylo bleu, grands
    caractères.
  \item \textbf{Pages 7 et 8 : espace classifiant d'un groupe discret.} Deux
    fiches de petit format, deux théorèmes, la construction de Borel.
  \item \textbf{Pages 9 et 10 : cartes et groupes cartographiques.} Pages
    d'esquisses au feutre, formules dispersées et dessins.
  \item \textbf{Pages 11 à 15 : $\pi$-ensembles et bitorseurs.} Papier fin,
    encre noire ; les pages 11 et 12 sont numérotées ① et ② de sa main. Le
    seul ensemble qui ait la forme d'une rédaction, avec Proposition et
    Corollaire.
\end{itemize}

Ce qui les rapproche est une question plutôt qu'un théorème : un groupe
classe les torseurs sous lui, et l'on cherche chaque fois l'objet qui porte
cette classification — le groupe $\varprojlim \mu_n$ de Kummer, l'espace $B\pi$,
le groupe cartographique, la catégorie $\mathrm{Ens}(\pi)$. Deux liens entre
les ensembles sont de cette lecture et non des pages : le groupe
$\pi_1(X, \Gamma, a)$ de la page 14 est le modèle combinatoire du
$\pi_1(X, \pi)$ de la page 8, et l'involution $\varepsilon$ des pages 11 et 14
est un exemple des extensions par $\mathbb{Z}/2$ de la page 13, dont la page 15
écrit la condition de cohérence.

Le titre de l'inventaire, « note Alexandre (bi tenseur…) », n'est éclairé par
aucune page ; « bitorseur », le mot des pages 12 à 15, en est une lecture
vraisemblable, rien de plus.

\textbf{Conventions, valables pour tout le dossier.} Les groupes agissent à
gauche sauf mention contraire ; $\mathrm{Ens}(\pi)$ est la catégorie des
$\pi$-ensembles à gauche. Un $(\pi', \pi)$-ensemble ${}_{\pi'}X_{\pi}$ porte une
action de $\pi'$ à gauche et une action de $\pi$ à droite qui commutent ; le
\emph{produit contracté} $X \times_{\pi} E$ d'un $\pi$-ensemble à droite $X$ et
d'un $\pi$-ensemble à gauche $E$ est le quotient de $X \times E$ par
$(x\lambda, e) \sim (x, \lambda e)$. Un $(\pi', \pi)$-\emph{bitorseur} est un
$(\pi', \pi)$-ensemble sur lequel chacune des deux actions est libre et
transitive ; $\mathrm{Bit}(\pi', \pi)$ désigne le groupoïde qu'ils forment. On
note $\mathrm{Int}(\pi)$ les automorphismes intérieurs,
$\mathrm{Isomext}(\pi', \pi) = \mathrm{Isom}(\pi', \pi)/\mathrm{Int}(\pi)$ et
$\mathrm{Out}(\pi) = \mathrm{Isomext}(\pi, \pi)$. Pour un espace $X$,
$H^1(X, \pi)$ est l'ensemble des classes d'isomorphisme de revêtements
principaux de groupe $\pi$.

\pagerange{1}{5}
\section*{Théorie de Kummer (pages 1, 3 et 5)}

\pagerange{1}{1}
\subsection*{Le corps des fractions d'un anneau de valuation discrète}

Soit $V$ un anneau de valuation discrète complet\footnote{La page écrit
« a.v.d. noethérien ; complet », lecture incertaine de l'abréviation ; un
anneau de valuation discrète est de toute façon noethérien. Ce qui sert
ci-dessous est que $V$ soit hensélien, ce qu'entraîne la complétude.}, de
corps résiduel $k$ et de corps des fractions $K$ ; soit $\varpi$ une
uniformisante et $\overline{K}$ une clôture algébrique de $K$.

\textbf{Cas où $k$ est algébriquement clos de caractéristique $0$.} Alors
\[
K^{*}/K^{*n} \simeq \mathbb{Z}/n\mathbb{Z}, \qquad x \longmapsto v(x) \bmod n,
\]
où $v$ est la valuation. En effet $K^{*} = \varpi^{\mathbb{Z}} \times V^{*}$, et
tout élément de $V^{*}$ est une puissance $n$-ième : sa réduction dans $k^{*}$
en est une, $k$ étant algébriquement clos, et l'on relève la racine par le
lemme de Hensel, puisque $n$ est inversible dans $k$\footnote{La page donne
l'isomorphisme « par la valuation » sans justification ; celle-ci est
ajoutée.}. Le même argument montre que $K$ contient les racines $n$-ièmes de
l'unité, et $\mu_n(K) = \mu_n(k)$. La théorie de Kummer identifie alors les
extensions abéliennes d'exposant divisant $n$ aux sous-groupes de
$K^{*}/K^{*n}$ : il y en a une seule de degré $n$, $K(\varpi^{1/n})$, et elle est
cyclique. Comme toute extension finie de $K$ est totalement et modérément
ramifiée — le corps résiduel est algébriquement clos et de caractéristique
nulle —, elle est de cette forme : \emph{toutes les extensions finies de $K$
sont cycliques, donc kummériennes}\footnote{« donc » est une lecture
incertaine ; l'argument est celui qu'on vient de dire.}.

Le groupe de Galois absolu se décrit alors canoniquement :
\[
\mathrm{Gal}(\overline{K}/K) \xrightarrow{\ \sim\ } T(k) := \varprojlim_n \mu_n(k),
\qquad \sigma \longmapsto \bigl(\sigma(\varpi^{1/n})/\varpi^{1/n}\bigr)_n ,
\]
et $T(k)$, le « module de Tate » du groupe multiplicatif de $k$, est isomorphe
à $\widehat{\mathbb{Z}}$, mais non canoniquement\footnote{La page note ce
groupe $T_\infty(k)$ et écrit $\Gamma(\overline{K}, K)$ pour le groupe de Galois.
L'isomorphisme ne dépend ni du choix des racines, les $\mu_n$ étant dans $K$, ni
de celui de l'uniformisante : si $\varpi' = u\varpi$ avec $u \in V^{*}$, les
racines $n$-ièmes de $u$ sont déjà dans $K$. On l'écrit aujourd'hui
$\widehat{\mathbb{Z}}(1)$.}.

\textbf{Cas où $k$ est de caractéristique $p > 0$.} On ne suppose plus $V$
complet, mais seulement hensélien\footnote{La page dit « si $k$ car $\neq 0$,
$V$ pas nécessairement complet » ; le signe « $\neq$ » est une lecture, et la
page ne dit pas quelle hypothèse remplace la complétude. Sans hensélianité,
les extensions de $K$ ne correspondent plus à celles du complété, et c'est le
groupe de décomposition qui joue le rôle de $\mathrm{Gal}(\overline{K}/K)$.}.
Les extensions de degré premier à $p$ du même type donnent le quotient
\emph{modéré}, et l'on a une suite exacte
\[
1 \to \varprojlim_{(n, p) = 1} \mu_n(\overline{K}) \longrightarrow \pi_1^{t}(K)
\longrightarrow \pi_1(k) \to 1, \qquad \pi_1(k) = \pi_1(\mathrm{Spec}\, V),
\]
où $\pi_1^{t}(K)$ est le groupe de Galois de l'extension modérément ramifiée
maximale, $\pi_1(k) = \mathrm{Gal}(k^{s}/k)$, et le noyau est l'inertie modérée,
isomorphe à $\prod_{\ell \neq p} \mathbb{Z}_\ell(1)$\footnote{La page écrit la
suite avec $\pi_1(K)$ au milieu (page 3) et, page 1, une flèche
$\pi_1(K)^{\text{premier à } p} \to \varprojlim_{(n,p)=1} \mu_n(k)$. Pour le
groupe $\pi_1(K)$ entier la suite n'est pas exacte à gauche : le noyau de
$\pi_1(K) \to \pi_1(k)$ est l'inertie tout entière, extension de l'inertie modérée
par un pro-$p$-groupe, l'inertie sauvage. L'égalité
$\pi_1(k) = \pi_1(\mathrm{Spec}\, V)$ est une propriété des anneaux henséliens. Le
groupe fondamental modéré est l'objet du livre de Grothendieck et Murre sur le
voisinage formel d'un diviseur à croisements normaux (1971).}.

\pagerange{3}{5}
\subsection*{Sur un schéma quelconque : la suite de Kummer}

Soit $S$ un schéma et $n$ un entier inversible sur $S$. Toute unité
$a \in \Gamma(S, \mathcal{O}_S^{*})$ définit un revêtement
\[
S_a = \mathrm{Spec}\bigl(\mathcal{O}_S[T]/(T^n - a)\bigr) \longrightarrow S,
\]
étale, torseur sous $\mu_n$ qui agit par $T \mapsto \zeta T$ ; sur un ouvert
affine $\mathrm{Spec}\, A$, c'est $A[T]/(T^n - a)$. Si $S$ est connexe et si
$\mu_n$ y est constant, un tel torseur, pointé au-dessus d'un point géométrique
$\bar{s}$, équivaut à un homomorphisme $\pi_1(S, \bar{s}) \to \mu_n(S)$ — le
caractère de Kummer de $a$ ; une famille compatible de racines de $a$ donne
une flèche $\pi_1(S, \bar{s}) \to \varprojlim_n \mu_n(S)$, qui généralise celle de
la page 1, où $a = \varpi$\footnote{La page 3 écrit
$\pi_1(S, \bar{s}) \to \varprojlim \mu_n(S)$ sans dire ce qui définit la flèche,
et place sous la limite une condition lue « $n$ premier à $s$ ». La lecture par
le caractère de Kummer d'une unité, et la condition « $n$ inversible sur $S$ »,
sont de cette édition.}.

Le classement de tous ces torseurs est la suite exacte longue associée à la
\emph{suite de Kummer} $1 \to \mu_n \to \mathbb{G}_m \xrightarrow{\,n\,} \mathbb{G}_m \to 1$,
exacte pour la topologie étale quand $n$ est inversible sur $S$, et pour la
topologie fppf en général\footnote{La page ne précise ni la topologie ni
l'hypothèse sur $n$. Elle note $\mathrm{Rev}(S, \mu_n)$ ce qui est ici
$H^1(S, \mu_n)$.} :
\[
0 \to \mu_n(S) \to \mathbb{G}_m(S) \xrightarrow{\,n\,} \mathbb{G}_m(S)
\xrightarrow{\ \partial\ } H^1(S, \mu_n) \to \mathrm{Pic}(S)
\xrightarrow{\,\otimes n\,} \mathrm{Pic}(S),
\]
en utilisant $H^0(S, \mathbb{G}_m) = \mathbb{G}_m(S)$ et
$H^1(S, \mathbb{G}_m) = \mathrm{Pic}(S)$\footnote{La première ligne de la page 5
omet le second $H^0(S, \mathbb{G}_m)$ et la puissance $n$-ième qui y mène ; la
seconde les rétablit.}. Le cobord $\partial$ est $a \mapsto [S_a]$, et
$H^1(S, \mu_n) \to \mathrm{Pic}(S)$ envoie un $\mu_n$-torseur sur le fibré en droites
qu'il induit par $\mu_n \subset \mathbb{G}_m$, dont la puissance $n$-ième est
trivialisée. D'où la suite exacte courte de la page :
\[
0 \to \mathbb{G}_m(S)/\mathbb{G}_m(S)^{n} \longrightarrow H^1(S, \mu_n)
\longrightarrow {}_{n}\mathrm{Pic}(S) \to 0,
\]
où ${}_{n}\mathrm{Pic}(S)$ est le sous-groupe de $n$-torsion\footnote{La page écrit
le premier terme $\mathbb{G}_m(S)_n$, et le dernier, surchargé, se lit
${}_n\mathrm{Pic}(S)$.} : un $\mu_n$-torseur est donné par une unité, à puissance
$n$-ième près, dès que le fibré en droites qu'il définit est trivial.

\pagerange{7}{8}
\section*{Espace classifiant d'un groupe discret (pages 7 et 8)}

Soit $\pi$ un groupe discret et $E \to B$ un revêtement principal de groupe
$\pi$. Pour tout espace $X$ — disons un CW-complexe\footnote{La page ne
restreint pas la classe des espaces. L'invariance par homotopie de l'image
réciproque, qui rend $h_X$ bien défini, demande $X$ paracompact ; l'équivalence
ci-dessous se démontre pour $X$ parcourant les CW-complexes.} — on dispose de
l'application, fonctorielle en $X$,
\[
h_X : [X, B] \longrightarrow H^1(X, \pi), \qquad f \longmapsto [f^{*}E],
\]
où $[X, B]$ désigne les classes d'homotopie d'applications\footnote{La page
écrit $\mathrm{hot}(X, B)$.}.

\textbf{Théorème 1.} Les conditions suivantes sont équivalentes :
\begin{enumerate}
  \item $h$ est un isomorphisme de foncteurs, i.e. $h_X$ est bijectif pour tout
    $X$ ;
  \item $E$ est faiblement contractile\footnote{La fiche s'arrête sur
    « 2) $E$ est contractile », coupée au bas ; elle portait peut-être d'autres
    conditions. « Contractile » est exact quand $B$ est un CW-complexe (théorème
    de Whitehead) ; « faiblement contractile » en général. La page corrige
    « injectif » en « bijectif » au stylo bleu.}.
\end{enumerate}
On dit alors que $(E, B, \pi)$ est un revêtement universel de groupe $\pi$, et
que $B$ est un \emph{espace classifiant} pour $\pi$.

L'équivalence tient à ceci. Appliquée à un point, la condition 1 dit que $B$ est
connexe ; appliquée aux sphères $S^k$, $k \geq 2$, pour lesquelles
$H^1(S^k, \pi)$ est réduit à un point, elle dit que $\pi_k(B) = 0$ ; appliquée à
$S^1$ et à $B$ lui-même, que $E$ est connexe et que $\pi_1(B) \simeq \pi$. Donc $E$
est le revêtement universel d'un espace $B$ dont tous les groupes d'homotopie
supérieurs sont nuls : c'est la condition 2. La réciproque est la théorie de
l'obstruction. En termes d'aujourd'hui, $B$ est un espace d'Eilenberg–MacLane
$K(\pi, 1)$, et $H^1(X, \pi) \simeq \mathrm{Hom}(\pi_1(X), \pi)/\text{conjugaison}$ pour
$X$ connexe\footnote{Cette justification est de l'édition ; la page énonce sans
démontrer.}.

\textbf{Théorème 2.} Pour tout groupe discret $\pi$, il existe un revêtement
universel de groupe $\pi$, et le couple formé de l'espace classifiant $B_\pi$ et de
la classe de $E_\pi$ est défini à isomorphisme unique près dans la catégorie
homotopique\footnote{La page dit « un espace classifiant (défini à isom. unique
près dans la cat. homotopique…) ». L'unicité à isomorphisme \emph{unique} vaut
pour le couple, qui représente le foncteur $H^1(-, \pi)$ ; l'espace $B_\pi$ seul
n'est unique qu'à isomorphisme près, ses auto-équivalences réalisant
$\mathrm{Out}(\pi)$. L'existence est due à Milnor (1956) pour les groupes
topologiques, et pour un groupe discret se lit sur la réalisation géométrique du
nerf de $\pi$.}.

\textbf{La construction de Borel.} Soit $X$ un espace sur lequel $\pi$
agit\footnote{Après « opère », la page porte quelques signes illisibles ; aucune
hypothèse sur l'action n'est nécessaire à ce qui suit.}. On pose
\[
X_{h\pi} = E_\pi \times_\pi X = (E_\pi \times X)/\pi ,
\]
qui se projette sur $B_\pi = E_\pi/\pi$ et sur $X/\pi$\footnote{Le but de la seconde
projection est surchargé sur la page ; $X/\pi$ est ce que la formule
impose.}. La première projection est un fibré de fibre $X$ ; la seconde est
une équivalence d'homotopie quand l'action est libre et proprement
discontinue et $X$ paracompact, $E_\pi$ étant contractile. La suite exacte d'homotopie de la
fibration $X \to X_{h\pi} \to B_\pi$ donne, puisque $\pi_2(B_\pi) = 0$ :
\[
0 = \pi_2(B_\pi) \to \pi_1(X) \longrightarrow \pi_1(X, \pi) \longrightarrow \pi
\longrightarrow \pi_0(X),
\]
où $\pi_1(X, \pi) := \pi_1(X_{h\pi})$, et le dernier terme est un ensemble pointé.
Si $X$ est connexe — la page écrit $\pi_0(X) = 1$ —, c'est une extension
\[
1 \to \pi_1(X) \to \pi_1(X, \pi) \to \pi \to 1 ,
\]
celle qu'on appelle aujourd'hui le groupe fondamental équivariant, ou le groupe
fondamental orbifolde quand l'action est propre. La page 14 en donnera, sans le
dire, une description combinatoire.

\pagerange{9}{10}
\section*{Cartes et groupes cartographiques (pages 9 et 10)}

Ces deux pages sont des esquisses : formules dispersées, flèches et dessins.
On en tire ce qui s'énonce, sans prétendre restituer un ordre qu'elles n'ont
pas.

\pagerange{9}{9}
\subsection*{Le groupe cartographique et les cartes de type $(p, q)$}

Une \emph{carte} est un graphe $K$ tracé sur une surface $X$, découpant la
surface en faces ; les pages l'écrivent $X \supset K \supset S$, surface,
graphe, ensemble des sommets. Un \emph{drapeau} est un triplet
(sommet, arête, face) incidents ; le groupe
\[
\Gamma = \langle \sigma_0, \sigma_1, \sigma_2 \mid \sigma_0^2 = \sigma_1^2 = \sigma_2^2 = (\sigma_0\sigma_2)^2 = 1 \rangle
\]
agit sur les drapeaux, $\sigma_i$ changeant la seule composante de dimension $i$.
Une carte connexe pointée par un drapeau $\tilde{x}$ — au sens large, surfaces à
bord, non orientables et cartes infinies comprises — équivaut à un sous-groupe
$\pi \subset \Gamma$, le stabilisateur de $\tilde{x}$ ; une carte connexe
non pointée, à une classe de conjugaison de sous-groupes ; et le groupe des
automorphismes de la carte est $\mathrm{Norm}(\pi)/\pi$. La catégorie des cartes
s'identifie ainsi à celle des $\Gamma$-ensembles, et le revêtement universel
$\widetilde{X} \supset \widetilde{K} \supset \widetilde{S}$ en est l'objet
correspondant à $\pi = 1$\footnote{Les pages ne définissent ni $\Gamma$ ni la
catégorie ; la présentation ci-dessus est celle que Grothendieck publiera dans
l'\emph{Esquisse d'un programme} (1984), sous le nom de groupe cartographique.
Les pages écrivent seulement
$\tilde{x} \in (\widetilde{X} \supset \widetilde{K} \supset \widetilde{S}) \simeq C_{p,q}$,
$\Gamma_{p,q} \supset \pi$ et $\mathrm{Norm}(\pi)/\pi$.}.

Les éléments pairs de $\Gamma$ forment le groupe cartographique \emph{orienté}
$\Gamma^{0}$, qui agit sur les cartes orientées ; il est engendré par les deux
rotations $\sigma_1\sigma_2$ (autour d'un sommet) et $\sigma_0\sigma_1$ (autour d'une
face) et par le retournement d'arête $\sigma_2\sigma_0$, d'ordre $2$, dont le produit
est trivial ; d'où $\Gamma^{0} \simeq \mathbb{Z} * \mathbb{Z}/2$.

Pour $p, q \in \{1, 2, \ldots, \infty\}$, le groupe $\Gamma_{p,q}$ s'obtient en
imposant la relation $\rho^{p} = 1$ à l'une des deux rotations et $\rho'^{\,q} = 1$ à
l'autre ($p$ ou $q$ infini : pas de relation) ; ses
ensembles à action forment la catégorie $C_{p,q}$ des cartes dont les sommets,
respectivement les faces, sont de valence divisant ces entiers\footnote{La page
ne dit pas laquelle des deux rotations est bornée par $p$. Elle écrit
$\rho_s^3 = 1$ dans une présentation de $\Gamma_{\infty,3}$, $\rho_s$ étant
vraisemblablement la rotation autour d'un sommet, et, à part, $\rho_s^{\,p}$
souligné : les deux mentions ne s'accordent qu'avec des conventions opposées.
On n'énonce ici que ce qui ne dépend pas de ce choix.}. En particulier
$\Gamma_{\infty,\infty} = \Gamma$, qui se surjecte sur tous les autres :
$\Gamma_{\infty,\infty} \to \Gamma_{\infty,3}$, $\Gamma_{\infty,\infty} \to \Gamma_{p,q}$.

\textbf{Divisibilité.} Si $p' \mid p$ et $q' \mid q$, on a une surjection
$\Gamma_{p,q} \to \Gamma_{p',q'}$, de noyau le sous-groupe distingué $R$ engendré par
les relations supplémentaires. Un $\Gamma_{p',q'}$-ensemble est un
$\Gamma_{p,q}$-ensemble sur lequel $R$ agit trivialement : $C_{p',q'}$ est une
sous-catégorie pleine de $C_{p,q}$, et le foncteur $C_{p,q} \to C_{p',q'}$,
$E \mapsto E/R$, en est l'adjoint à gauche\footnote{La page écrit
« $C_{p',q'} = C_{p,q}/R$ » et $R_{(p,p',q,q')}$. L'égalité vaut pour les
objets universels — la carte universelle de type $(p', q')$ est le quotient par
$R$ de celle de type $(p, q)$ — ; pour les catégories, c'est un plongement plein
muni d'un adjoint à gauche.}. Les cas extrêmes $C_{1,1}$, $C_{1,2}$, $C_{2,1}$
figurent en bas de page.

\textbf{Triangulations et groupe modulaire.} Quand l'une des rotations est
d'ordre $3$ et l'autre libre, le groupe orienté est
\[
\Gamma^{0}_{\infty,3} \simeq \mathbb{Z}/2 * \mathbb{Z}/3 \simeq PSL(2, \mathbb{Z}) = SL(2, \mathbb{Z})/\{\pm 1\},
\]
et le groupe entier $\Gamma_{\infty,3}$ est le groupe modulaire étendu
$PGL(2, \mathbb{Z})$, groupe de réflexions du triangle hyperbolique d'angles
$\pi/2$, $\pi/3$, $0$. La réduction modulo $2$ donne une suite exacte
\[
1 \to \Gamma(2)/\{\pm 1\} \longrightarrow PSL(2, \mathbb{Z}) \longrightarrow SL(2, \mathbb{F}_2) \simeq \mathfrak{S}_3 \to 1,
\]
dont le noyau est libre de rang $2$ et s'identifie au groupe fondamental de la
sphère privée de trois points\footnote{La page écrit la suite verticalement,
avec un premier terme qui ressemble à $\Gamma_1$. Le noyau de
$PSL(2,\mathbb{Z}) \to \mathfrak{S}_3$ est $\Gamma(2)/\{\pm 1\}$ : c'est la seule lecture
compatible avec l'exactitude, et elle est de l'édition.}. La page le rattache
à la géométrie : un réseau de $\mathbb{C}$ de base $(\omega_1, \omega_2)$ avec
$\mathrm{Im}(\omega_2/\omega_1) > 0$, le module $\tau = \omega_2/\omega_1$, et, dessinée
à côté, la droite des modules — une sphère marquée de deux points exceptionnels
et d'une pointe —, dont $PSL(2, \mathbb{Z})$ est le groupe fondamental au sens des
orbifoldes\footnote{Les deux croix du dessin sont marquées $0$ et d'une valeur
lue $2^{3}\cdot 3^{4}$. S'il s'agit des valeurs exceptionnelles de l'invariant
modulaire, la seconde est $j = 1728 = 2^{6}\cdot 3^{3}$. Ce lien entre cartes
triangulées et courbes modulaires est celui que développera l'\emph{Esquisse
d'un programme} ; le théorème de Belyi (1979), qui en fait une théorie
arithmétique, était paru.}.

\pagerange{10}{10}
\subsection*{Cartes trouées}

La page 10 esquisse une définition topologique des cartes. Une \emph{carte
trouée} est un couple $(X \supset K)$ formé d'une surface topologique $X$ et
d'une partie fermée $K$ de dimension $1$\footnote{« sous-variété fermée de dim
$1$ » est une lecture incertaine ; la page s'arrête après « tels que ».},
tel que les composantes de $\text{Déc}(X, K)$ — la surface découpée le long de
$K$ — soient homéomorphes à des faces standard $D_n^{*}$, $1 \leq n \leq \infty$.

La face standard se construit à partir d'un polygone : soit $\Pi_n$ le polygone
convexe régulier à $n$ côtés, $\Pi_n = (S_n, A_n, F_n)$ ses sommets, arêtes et
face\footnote{« pol. conv. », « côtés » et la troisième lettre sont des lectures
incertaines.}, et $|\Pi_n|$ son bord, qui contient $S_n$. Alors
$|\Pi_n| \times \left]0,1\right]$ est un disque fermé privé de son centre — le cône
sur le bord, sommet ôté —, et
\[
D_n^{*} = |\Pi_n| \times \left]0,1\right] \smallsetminus S_n ,
\]
les sommets étant retirés du bord\footnote{L'interprétation de $|\Pi_n|$ comme
bord du polygone et de $|\Pi_n| \times \left]0,1\right]$ comme disque épointé est
de l'édition ; la page ne précise pas à quel niveau les sommets sont ôtés. Elle
est confirmée par les dessins : un pentagone découpé en secteurs depuis son
centre, et, pour $n = \infty$, une droite subdivisée prolongée à l'infini des
deux côtés, dont le produit par $\left]0,1\right]$ est un demi-plan.}. Les faces
d'une carte trouée sont donc épointées en leur centre : les composantes de
$\text{Déc}(X, K)$ ne sont pas compactes, même pour une carte finie. Les dessins de la page — une
face circulaire tournée de $\pm 2\pi/n$, une carte pentagonale aux sommets
cerclés et arêtes pendantes, les générateurs $\sigma_0, \sigma_1, \sigma_2$ — restent
sans texte.

\pagerange{11}{15}
\section*{$\pi$-ensembles et bitorseurs (pages 11 à 15)}

\pagerange{11}{11}
\subsection*{Deux rappels : indice des sous-groupes libres, sphère privée de trois points}

Si $L$ est un groupe libre de rang $d$ et $L' \subset L$ un sous-groupe d'indice
$n$, alors $L'$ est libre de rang $d'$ avec
\[
d' - 1 = n(d - 1)
\]
(formule de Schreier ; par exemple $d = 2$, $n = 2$ donne $d' = 3$). Topologiquement :
si $L = \pi_1(X)$ pour un bouquet $X$ de $d$ cercles, $L'$ est le groupe
fondamental d'un revêtement $X'$ à $n$ feuillets, et la caractéristique
d'Euler, $\chi(X) = 1 - d$, est multipliée par $n$.

La sphère privée des trois points $0, 1, \infty$ a pour groupe fondamental
$\pi^{+} = \langle \rho_0, \rho_1, \rho_\infty \mid \rho_\infty \rho_1 \rho_0 = 1 \rangle$, libre
de rang $2$. La conjugaison complexe, qui fixe l'axe réel et donc les trois
points, y induit l'automorphisme involutif
\[
\rho_0 \mapsto \rho_0^{-1}, \qquad \rho_\infty \mapsto \rho_\infty^{-1}, \qquad
\rho_1 \mapsto \rho_\infty \rho_0 = \rho_0^{-1}\,\rho_1^{-1}\,\rho_0 ,
\]
défini sur les générateurs libres $\rho_0, \rho_\infty$\footnote{La page ne donne pas
la relation. La formule pour $\rho_1$ n'est compatible qu'avec
$\rho_\infty \rho_1 \rho_0 = 1$ (ou sa permutation circulaire
$\rho_1\rho_0\rho_\infty = 1$) : avec $\rho_0\rho_1\rho_\infty = 1$ il faudrait
$\rho_1 \mapsto \rho_0 \rho_\infty$. Le lien avec la conjugaison complexe, suggéré
par le triangle hachuré de sommets $0$, $1$, $\infty$ dessiné à côté, est de
l'édition.}. On le retrouvera page 14.

\pagerange{11}{12}
\subsection*{Foncteurs sur les $\pi$-ensembles qui commutent aux limites inductives}

Soit $\mathcal{C}$ une catégorie cocomplète\footnote{La page écrit $\mathcal{C}$ reliée à
un mot lu « avec $\varinjlim$ » ; l'hypothèse est nécessaire pour que les produits
contractés ci-dessous existent.}. On a une équivalence entre
\begin{itemize}
  \item les foncteurs $F : \mathrm{Ens}(\pi) \to \mathcal{C}$ qui commutent aux limites
    inductives,
  \item les objets $X$ de $\mathcal{C}$ munis d'une action à droite de $\pi$,
\end{itemize}
donnée par $F \mapsto F(\pi_s)$, où $\pi_s$ est $\pi$ agissant sur lui-même à
gauche, et dont l'action de $\pi$ à droite par translations fait de $F(\pi_s)$ un
$\pi$-objet à droite ; en sens inverse, $X \mapsto F_X$ avec
\[
F_X(E) = X \times_\pi E, \qquad
\mathrm{Hom}_{\mathcal{C}}(X \times_\pi E, Y) \simeq \mathrm{Hom}_\pi\bigl(E, \mathrm{Hom}_{\mathcal{C}}(X, Y)\bigr).
\]
La formule d'adjonction dit que $F_X$ a un adjoint à droite, donc commute aux
limites inductives ; qu'on retrouve ainsi tout $F$ vient de ce que tout
$\pi$-ensemble est somme d'orbites,
\[
E \simeq \coprod_{i \in I} \pi/H_i, \qquad
X \times_\pi E \simeq \coprod_{i} X \times_\pi (\pi/H_i) = \coprod_i X/H_i ,
\]
où $X/H_i$ est le quotient dans $\mathcal{C}$. En termes d'aujourd'hui,
$\mathrm{Ens}(\pi)$ est la catégorie des préfaisceaux sur le groupe vu comme
catégorie à un objet, sa \emph{complétion par limites inductives libre}, et $F_X$
est l'extension de Kan à gauche de $X$\footnote{La page écrit $X$ « $\pi$-objet (à
droite) de $\mathcal{C}$ », « à droite » ajouté en interligne. Elle esquisse aussi
un foncteur $G : \mathcal{C} \to \mathcal{C}'$ commutant aux limites inductives, sans
suite : le composé $G \circ F_X$ correspond à $G(X)$.}.

\textbf{Exemple : $\mathcal{C} = \mathrm{Ens}(\pi')$.} Un $\pi'$-ensemble muni d'une action à
droite de $\pi$ est un $(\pi', \pi)$-ensemble, et l'on obtient une équivalence
\[
\mathrm{Hom}_{!}\bigl(\mathrm{Ens}(\pi), \mathrm{Ens}(\pi')\bigr) \simeq {}_{\pi'}\mathrm{Ens}_{\pi},
\qquad F_X \longleftrightarrow {}_{\pi'}X_{\pi},
\]
où $\mathrm{Hom}_{!}$ désigne les foncteurs qui commutent aux limites
inductives\footnote{C'est la définition que la page accole au « ! ».}. Un
$(\pi', \pi)$-ensemble est un foncteur $B\pi' \times (B\pi)^{\mathrm{op}} \to \mathrm{Ens}$, où
$B\pi$ désigne le groupe vu comme catégorie à un objet : c'est un
\emph{profoncteur} de $B\pi$ vers $B\pi'$, et la composition des profoncteurs est
le produit contracté (page 15).

\textbf{Proposition.} Pour que $F_X$ soit exact à gauche, il faut et il suffit
que $X$ soit un $\pi$-torseur à droite, i.e. que l'action à droite de $\pi$ sur $X$
soit libre et transitive.

En effet, $F_X$ doit préserver l'objet final, donc $X/\pi$ est un point :
$X \simeq H\backslash\pi$ pour un sous-groupe $H$, et alors $F_X(E) = H\backslash E$,
l'ensemble des orbites de $H$. Il doit aussi préserver les produits ; or dans
$H\backslash(\pi_s \times \pi_s)$, les classes de $(1, 1)$ et de $(1, h)$, $h \in H$, ont
même image dans $H\backslash\pi \times H\backslash\pi$ et ne sont égales que si
$h = 1$ : donc $H = 1$. Réciproquement, si $X$ est un torseur, le choix d'un
point identifie $F_X(E)$ à $E$ lui-même, et $F_X$ préserve les limites, qui se
calculent sur les ensembles sous-jacents\footnote{Démonstration de l'édition ; la page énonce seulement.
« Exact à gauche » s'entend : préserve les limites projectives finies.}.

Pour un $\pi$-torseur à droite $X$, on note
\[
{}^{X}\pi := X \times_\pi \pi_{\mathrm{int}} = \mathrm{Aut}_\pi(X),
\]
la \emph{forme tordue} de $\pi$ par $X$, où $\pi_{\mathrm{int}}$ est $\pi$ agissant sur
lui-même par conjugaison ; le choix d'un point $x \in X$ l'identifie à $\pi$, et
deux choix diffèrent d'un automorphisme intérieur. Une structure de
$(\pi', \pi)$-ensemble sur $X$ est alors un homomorphisme $\pi' \to \mathrm{Aut}_\pi(X)$.

\textbf{Corollaire.} Pour que $F_X$ soit une équivalence, il faut et il suffit
que $\pi' \to \mathrm{Aut}_\pi(X)$ soit un isomorphisme, i.e. que $X$ soit un
$(\pi', \pi)$-bitorseur. D'où une équivalence de groupoïdes
\[
\text{Équiv}\bigl(\mathrm{Ens}(\pi), \mathrm{Ens}(\pi')\bigr) \simeq \mathrm{Bit}(\pi', \pi),
\]
qui envoie un quasi-inverse $u^{[-1]}$ sur le bitorseur opposé $X^{-1}$ — le même
ensemble, $\pi$ agissant à gauche et $\pi'$ à droite par les inverses\footnote{La
page dessine le carré, avec $u \mapsto u^{[-1]}$ à gauche ; l'étiquette de la
flèche de droite est illisible, et « bitorseur opposé » est ce que la
commutativité du carré impose.}.

\pagerange{13}{13}
\subsection*{Bitorseurs involutifs et extensions par $\mathbb{Z}/2$}

Un $(\pi, \pi)$-bitorseur $X$ est trivial — isomorphe au bitorseur $\pi$ — si et
seulement s'il contient un élément $e$ tel que $\lambda e = e \lambda$ pour tout
$\lambda \in \pi$. Un bitorseur \emph{involutif} est un $(\pi, \pi)$-bitorseur $X$ muni
d'un isomorphisme $\mu : X \times_\pi X \xrightarrow{\sim} \pi$, associatif au sens où les
deux isomorphismes $X \times_\pi X \times_\pi X \to X$ qu'il induit coïncident\footnote{La
page écrit $X \circ X \xrightarrow{\sim} 1$ sans condition ; la condition
d'associativité est celle que la page 15 écrit pour les foncteurs,
$F\mu = \mu F$.}. La page affirme des équivalences de groupoïdes
\[
\mathrm{Bit.invol}(\pi) \simeq \mathrm{Invol}\bigl(\mathrm{Ens}(\pi)\bigr) \simeq \mathrm{Ext}(\mathbb{Z}/2, \pi).
\]
La première est le Corollaire des pages 11 et 12, appliqué à $\pi' = \pi$ et
complété par $\mu$ : une involution de $\mathrm{Ens}(\pi)$ est un foncteur $F$ avec
$\mu : F^2 \xrightarrow{\sim} \mathrm{id}$ et $F\mu = \mu F$ (page 15). La seconde
associe à une extension $1 \to \pi \to \widetilde{\pi} \xrightarrow{p} \mathbb{Z}/2 \to 1$ la
classe non neutre $X = p^{-1}(1)$, qui est un $(\pi, \pi)$-bitorseur par
multiplication à gauche et à droite, et que la multiplication de $\widetilde{\pi}$
munit de $\mu : X \times_\pi X \to \pi$ ; inversement $\widetilde{\pi} = \pi \sqcup X$, la
loi étant donnée par les actions et par $\mu$\footnote{Justification de
l'édition ; la page énonce. Le classement des extensions par
$H^2(\mathbb{Z}/2, Z(\pi))$ au-dessus d'une action extérieure donnée est la théorie de
Schreier et d'Eilenberg–MacLane.}.

\textbf{Le groupe d'un polygone idéal.} Soit $I$ un ensemble fini à $n \geq 3$
éléments, qu'on pense comme les côtés d'un polygone convexe\footnote{« pol.
conv. fini » après un mot biffé, puis « cartes $n$-gonales
$I$-pondérées » : lectures incertaines. La géométrie hyperbolique qui suit est
de l'édition.}. Le produit libre
\[
\pi_I = (\mathbb{Z}/2)^{*I}
\]
est le groupe engendré par les réflexions par rapport aux côtés d'un polygone
hyperbolique idéal à $n$ côtés — deux côtés adjacents se rencontrant à l'infini,
aucune relation ne lie leurs réflexions. Le morphisme $\pi_I \to \mathbb{Z}/2$ qui envoie
chaque générateur sur $1$ a pour noyau le sous-groupe des déplacements, et
\[
1 \to \pi_I^{+} \to \pi_I \to \mathbb{Z}/2 \to 1, \qquad
\pi_I^{+} \simeq L_I/(\lambda_0 \lambda_1 \cdots \lambda_{n-1}),
\]
groupe libre de rang $n - 1$, groupe fondamental de la sphère privée de $n$
points — le polygone doublé le long de son bord\footnote{La page écrit
d'abord $\pi_I = L_I/(\lambda_0\lambda_1\cdots\lambda_{n-1} = 1)$ et le biffe, puis
$\pi_I = (\mathbb{Z}/2)^{*I}$ et $\pi_I^{+} \simeq L_I/\!-$ : la présentation biffée était
celle du sous-groupe orienté, non de $\pi_I$. Le rang $n - 1$ se vérifie par la
caractéristique d'Euler rationnelle, $\chi(\pi_I) = 1 - n/2$ et
$\chi(\pi_I^{+}) = 2 - n$. Un dessin de carte en quadrilatères aux sommets
numérotés accompagne la suite.}. C'est un exemple de l'extension par
$\mathbb{Z}/2$ qu'on vient de décrire : la réflexion par rapport à un côté est un
élément de la classe non neutre.

\pagerange{14}{14}
\subsection*{L'involution sur un groupoïde fondamental, et le groupe $\pi_1(X, \Gamma, a)$}

Pour $n = 4$, la page prend, semble-t-il, la sphère de Riemann privée de
quatre points de l'axe réel, $\pi$ engendré par les lacets $\lambda_0, \lambda_1, \lambda_2, \lambda_3$ avec
$\lambda_3\lambda_2\lambda_1\lambda_0 = 1$\footnote{Après une première écriture biffée
$\lambda_0\lambda_1\lambda_2\lambda_3 = 1$ ; comme page 11, l'ordre du produit
est celui qui rend les formules de l'involution cohérentes.}, et
l'involution $\varepsilon$ induite par la conjugaison complexe :
\[
\varepsilon^2 = 1, \qquad \varepsilon\lambda_0\varepsilon^{-1} = \lambda_0^{-1}, \qquad
\varepsilon\lambda_1\varepsilon^{-1} = \lambda_1^{-1},
\]
et note à côté $\lambda_1^{-1}\lambda_2^{-1}\lambda_1$, vraisemblablement l'image de
$\lambda_2$, conjuguée de $\lambda_2^{-1}$ comme page 11\footnote{L'attribution de
cette expression à $\varepsilon\lambda_2\varepsilon^{-1}$, et la lecture des
quatre points comme réels, sont de l'édition : la page dessine une sphère
marquée de $i$, $-i$, $1$, d'un point à gauche et de $0$ au centre.}.
La difficulté que la page affronte est que la conjugaison ne fixe pas le point
base : elle échange $i$ et $-i$. On travaille donc dans le groupoïde
fondamental, où $\mathrm{Hom}(i, i)$ et $\mathrm{Hom}(-i, i)$ sont respectivement un
groupe et un torseur sous $\pi$, et l'automorphisme de $\pi$ s'obtient en
transportant par un chemin : $\lambda \mapsto \varepsilon \circ \bar\lambda \circ \varepsilon^{-1}$,
où $\bar\lambda$ est le lacet conjugué, pointé en $-i$\footnote{La page dessine les
lacets autour de $0$ et $1$ pointés en $i$ et $-i$, et trois chemins composables
dont les extrémités sont en partie surchargées ; leur lecture est
incertaine.}.

La seconde moitié de la page généralise. Soit $\Gamma$ un groupe agissant sur un
espace $X$ et $a \in X$. On pose
\[
\pi_1(X, \Gamma, a) = \{(\gamma, \ell) \mid \gamma \in \Gamma, \ \ell \in \mathrm{Hom}(\gamma a, a)\},
\]
les $\ell$ étant des classes de chemins, avec la loi
\[
(\gamma', \ell')(\gamma, \ell) = \bigl(\gamma'\gamma, \ \ell' \circ \gamma'(\ell)\bigr) :
\quad \gamma'\gamma a \xrightarrow{\ \gamma'(\ell)\ } \gamma' a \xrightarrow{\ \ell'\ } a .
\]
La page écrit la loi plus généralement pour des chemins
$\ell : \gamma a \to b$ et $\ell' : \gamma' b \to c$, ce qui en fait la composition d'un
groupoïde\footnote{Une première écriture, à gauche de la page, compose dans
l'autre ordre, $(\gamma\gamma', (\gamma\ell')\circ\ell)$, et les extrémités ne s'y
raccordent pas ; elle est en partie biffée. La loi retenue est celle que la page
écrit à droite.}. La projection $(\gamma, \ell) \mapsto \gamma$ est un morphisme vers
$\Gamma$, de noyau $\pi_1(X, a)$, et l'image est formée des $\gamma$ qui ne changent pas
la composante connexe par arcs de $a$. C'est exactement l'extension
$1 \to \pi_1(X) \to \pi_1(X, \pi) \to \pi$ de la page 8 : quand $\Gamma = \pi$ est
discret et $X$ assez régulier, $\pi_1(X, \Gamma, a) \simeq \pi_1(E_\pi \times_\pi X)$\footnote{Ce rapprochement est
de l'édition. Pour $\Gamma = \mathbb{Z}/2$ agissant par conjugaison complexe, on
retrouve l'extension de la page 13 ; pour une action propre, c'est le groupe
fondamental orbifolde du quotient.}.

\pagerange{15}{15}
\subsection*{Composition et classement des bitorseurs}

Le produit contracté compose les foncteurs : pour ${}_{\pi'}X_{\pi}$ et
${}_{\pi''}Y_{\pi'}$,
\[
\mathrm{Ens}(\pi) \xrightarrow{\ F_X\ } \mathrm{Ens}(\pi') \xrightarrow{\ F_Y\ } \mathrm{Ens}(\pi''),
\qquad F_Y \circ F_X \simeq F_{Y \times_{\pi'} X}
\]
— un isomorphisme canonique plutôt qu'une égalité, le produit contracté n'étant
associatif qu'à isomorphisme près\footnote{La page écrit l'égalité
$F_Y \circ F_X = F_{Y \wedge_{\pi'} X}$, avec la notation $\wedge$ pour le produit
contracté.}.

Comment classer les bitorseurs ? Si ${}_{\pi'}X_{\pi}$ est un bitorseur, on a
$\pi' \xrightarrow{\sim} {}^{X}\pi$, et ${}^{X}\pi \simeq \pi$ non canoniquement, mais
canoniquement à automorphisme intérieur de $\pi$ près : $X$ définit donc un
élément de $\mathrm{Isomext}(\pi', \pi)$\footnote{La page écrit en interligne « non
can. » puis « mais can. à $\mathrm{Isomext}({}^X\pi, \pi)$ près », lecture en partie
incertaine.}. On obtient ainsi
\[
\pi_0\bigl(\mathrm{Bit}(\pi', \pi)\bigr) \xrightarrow{\ \sim\ } \mathrm{Isomext}(\pi', \pi),
\]
et le même raisonnement pour les bitorseurs pointés, dont le point fixe
l'identification, donne le carré

\begin{tikzcd}
\pi_0\bigl(\mathrm{Bit}_{\mathrm{pt}}(\pi', \pi)\bigr) \arrow[r, "\sim"] \arrow[d] & \mathrm{Isom}(\pi', \pi) \arrow[d] \\
\pi_0\bigl(\mathrm{Bit}(\pi', \pi)\bigr) \arrow[r, "\sim"] & \mathrm{Isomext}(\pi', \pi)
\end{tikzcd}

où les flèches verticales sont surjectives, celle de droite étant le passage au
quotient par $\mathrm{Int}(\pi)$\footnote{La page écrit la flèche de droite en
oblique, annotée « surj. » et « passage au quotient par $\mathrm{Int}(\pi)$ »,
lecture incertaine ; « Classes de bit$(\pi', \pi)$ » est biffé au-dessus.}. Le
groupoïde $\mathrm{Bit}(\pi', \pi)$ n'est pas pour autant discret : le groupe des
automorphismes d'un bitorseur est le centre $Z(\pi)$. Pour $\pi' = \pi$,
$\mathrm{Bit}(\pi, \pi)$ est le $2$-groupe $\mathrm{AUT}(\pi)$, avec
$\pi_0 = \mathrm{Out}(\pi)$ et $\pi_1 = Z(\pi)$ — c'est le module croisé
$\pi \to \mathrm{Aut}(\pi)$ — et ce qui vient d'être dit est qu'il s'identifie aux
auto-équivalences de $\mathrm{Ens}(\pi)$\footnote{Ni le centre, ni le $2$-groupe, ni
les modules croisés (Whitehead, 1949) ne sont nommés par la page. Les
bitorseurs et leur rôle en cohomologie non abélienne sont exposés par Giraud
(1971).}.

\textbf{Involutions de catégories.} La page s'achève sur la définition qu'utilise
la page 13 : une involution d'une catégorie est un foncteur $F$ muni d'un
isomorphisme $\mu : F^2 \xrightarrow{\sim} \mathrm{id}$ tel que
\[
F\mu = \mu F : F^3 \longrightarrow F,
\]
i.e. tel que, pour tout objet $E$, les deux flèches
$F(\mu_E), \mu_{F(E)} : F(F(F(E))) \to F(E)$ coïncident. C'est exactement la condition
pour que $(F, \mu)$ définisse une action du groupe $\mathbb{Z}/2$ sur la catégorie à
isomorphisme cohérent près ; sur $\mathrm{Ens}(\pi)$, elle correspond à
l'associativité de $\mu$ pour un bitorseur involutif, et donc, par la page 13,
à une extension de $\mathbb{Z}/2$ par $\pi$.

\end{document}
